\scanpage{401}
\mspage{203}

arithmétisation de $\pi'_i$ --- ce qui
reviendrait à la bijectivité de
\[
\widehat{\widehat{\mathcal T}}/\mathcal N_a
\longrightarrow
\unsure{\widehat{\widehat{\mathcal T}}'}/\unsure{M_{a'}^{\sharp}|M_{a'}^{\natural}}
\simeq \widehat{\widehat{\mathcal T}}'/\mathcal N_{a'},
\]
qui \inserted{n'aurait} guère de raison d'être que si
on admettait
$\widehat{\widehat{\mathcal T}}(\pi)\xrightarrow{\sim}\unsure{\widehat{\widehat{\mathcal T}}}(\pi'_i)$,
qui me semble bien problématique.

\margin{\figure{[Deux traits verticaux parallèles bordent le paragraphe qui commence par « Cependant » et finit par $(g,\nu-1)$.]}}
Cependant, on trouve, par le foncteur
"bouchage des trous", une équivalence
entre la catégorie des \deleted{\unsure{courbes alg.}}
groupes profinis à lacets \uuline{extérieurs} \uline{arithmétisés},
\inserted{de type $(g,\nu)$,}
et des groupes profinis à lacets
(pas extérieurs !) \uline{arithmétisés}, de type
$(g,\nu-1)$.

On peut se proposer d'essayer de
préciser \deleted{\unsure{par}} le type des propriétés
qui vont caractériser $\mathcal N_{g,\nu}$ dans $\mathcal N'_{g,\nu}=$
$\operatorname{Norm}_{\widehat{\widehat{\mathcal T}}_{g,\nu}}
(\widehat{\mathcal T}_{g,\nu}^{+})$,
\deleted{Je \unsure{pense} qu'il}
\inserted{\deleted{\unsure{On aura}}\unsure{ou encore}}
$\unsure{\Gamma_{g,\nu}|\Pi_{g,\nu}}$
\deleted{\unsure{pour} \illegible{1} $\Gamma'$}
\unsure{dans} $\Gamma'_{g,\nu}=\mathcal N'_{g,\nu}/\widehat{\mathcal T}_{g,\nu}^{+}$.
On a des hom. naturels

\margin{400}

\scanpage{402}

\[
\tag{15}
\left|
\begin{aligned}
\Gamma'_{g,\nu}&\longrightarrow\operatorname{Autext}(\widehat{\mathcal T}_{g,\nu}^{+})\\
\Gamma'_{g,\nu}&\longrightarrow\operatorname{Aut\,ext}(\widehat{\mathcal T}_{g,\nu}^{+!})
\end{aligned}
\right.
\]
et je présume \unsure{que|qu'}\deleted{il s'agit} $\Gamma_{g,\nu}$ pourra
se décrire comme \unsure{l'image inverse} d'un
sous-groupe fermé convenable de l'un
ou l'autre des \unsure{deux} \unsure{mbs}, i.e. qu'on
peut le décrire en termes de propriétés
d'opérations extérieures sur $\widehat{\mathcal T}_{g,\nu}^{+}$ ou
sur $\widehat{\mathcal T}_{g,\nu}$. Les conditions \textcircled{e} en tous
cas sont bien de ce type (propriétés
de normaliser des \unsure{ss-groupes} \unsure{inscrits|associés}
\deleted{$\unsure{\mathcal T'_i}$} $\pi'_i$ de $\mathcal T_{g,\nu}^{+!}$ ...).
Bien sûr, on
pourrait poser des conditions sur des
actions extérieures (de $\widehat{\mathcal T}'{}_{g,\nu}^{+!}$, disons),
qui soient stables par \unsure{passages|passage} \inserted{\unsure{successifs.}} à des
$\widehat{\mathcal T}'{}_{g,\nu-1}^{+!}$, $\widehat{\mathcal T}'{}_{g,\unsure{\nu-2}}^{+!}$ etc ... et
\deleted{pour lesquelles}\inserted{qui soient}
engendrées par les conditions \textcircled{e}. Mais
il n'est pas dit du tout que cela

\margin{401}

\scanpage{403}
\mspage{204}

\unsure{suffise} à décrire les \deleted{$\unsure{\mathcal N_{g,\nu}}$}
$\Gamma_{g,\nu}\subset\Gamma'_{g,\nu}$,
ne serait-ce que parce que la condition
devient vide pour \deleted{$\unsure{\nu=0}$} le cas limite $\nu=0$
(si $g\geq2$; ou pour les cas $g=1$, $\nu=1$, ou $g=0$, $\nu=3$).

Il est possible qu'il faille faire
intervenir des propriétés des $\pi_1$ des
$\unsure{M_{g,\nu\,\overline{\mathbb Q}_a}}$, liés à la compactification.
Ce n'est guère que dans le cas de
$(g,\nu)=(0,3)$ qu'il ne faut pas \unsure{s'attendre}
\unsure{des tas !|du tout !} en guise de conditions.

Appelons "\unsure{courbe algébrique potentielle|courbes algébriques potentielles}"
(\unsure{sous-entendu}, sur une clôture alg. \unsure{pas|non} précisée de $\mathbb Q$,)
la donnée d'\unsure{une|un} \inserted{\unsure{groupe} \uline{\unsure{extprofini|catprofini}}}
$\pi$ à \unsure{lacets|lacet} \deleted{\unsure{profini}}\inserted{\deleted{ext}}
de type $(g,\nu)$, muni d'une arithmétisation
$a$, et d'un germe de relèvement
"admissible"
\margin{\deleted{\illegible{1}}}
\[
\tag{16}
\Gamma_{\unsure{a|\bullet}}^{\unsure{\natural}}
\longrightarrow\mathcal N_a.
\]
Elles forment une catégorie (pour les
iso, pour le moment) Si on se
donne un torseur $P$ sous $\Gamma_{g,\nu}$
(ce qui \inserted{\unsure{trivialmt|trivialemt}} revient à se donner une clôture

\margin{402}

\scanpage{404}
alg. $\overline{\mathbb Q}$ de $\mathbb Q$ \ldots) les courbes potentielles
de type $(g,\nu)$ \uline{relatives} à ce torseur
(\unsure{constant}, \unsure{correspondant|correspondent} à des courbes
alg. sur $\overline{\mathbb Q}$ \ldots) sont \inserted{alors} \unsure{munies} en
plus d'un \unsure{isom.} (``extérieur'')
\[
(17)\qquad \unsure{P_a|\mathbb P_a}\xrightarrow{\sim}P
\]
(définissant un isomorphisme
\[
(18)\qquad \unsure{\Gamma_a}\longrightarrow\Gamma_P
 =\operatorname{Aut}_{\unsure{\Gamma_{g,\nu}|\Pi_{g,\nu}}}(P).
\]
\margin{\unsure{Le cas où} $P$ est le torseur
\unsure{trivial, est celui|trivial, et celui} des \deleted{\illegible{1}} $\pi$ ext\illegible{1}
\unsure{munis} d'une \uline{\unsure{prédiscrétification}} \unsure{arbitraire}
(\unsure{d'où} une \deleted{\illegible{1}} \unsure{exacte|ayant}
$1\longrightarrow\unsure{\widehat{\mathcal T}_a}\longrightarrow
\unsure{\mathcal N_a}\longrightarrow\unsure{\Gamma_a|\Pi_a}\longrightarrow1$)}

Cette \unsure{fois-ci}, on s'attend à trouver
des \unsure{ens} d'iso. \uline{finis}, au lieu de
torseurs sous des groupes profinis
considérables.

On \illegible{1} les \deleted{\illegible{1}} \uline{Points} de $\pi$,
comme les classes de \inserted{$\pi$-}conjugaison
des \unsure{relèvts} ``admissibles'' de (16),
\unsure{ou}
\[
(19)\qquad\unsure{\Gamma_{\bullet}^{\natural}}\longrightarrow
\unsure{\mathcal M_a}\subset\widehat{\widehat{\mathcal T}}(\pi)
\]
i.e. un germe de vraie action de $\Gamma_{\bullet}$
sur $\pi$, avec
\[
(20)\qquad\pi^{\unsure{\Gamma_{\bullet}^{\natural}}}=1
\]
\figure{[L'égalité est encadrée par le haut et sur les deux côtés; l'encadrement est ouvert en bas.]}
(ceci pour \deleted{\illegible{1}}
\margin{403}

\scanpage{405}
\mspage{205}
l'admissibilité.

\uline{Question liminaire} La connaissance de
$\pi$, de $\underbrace{\widehat{\mathcal T}_a^+}_{\Sigma_a}
\subset\widehat{\widehat{\mathcal T}}(\pi)$, et \unsure{d'un} ss-groupe
$\Gamma'\subset\widehat{\widehat{\mathcal T}}(\pi)$
(normalisant $\Sigma_a$, \inserted{\unsure{tel que}
$\Gamma'\cap\unsure{\Sigma_a|\Sigma_{\bullet}}=(1)$})
\unsure{permettrait-il} de retrouver l'\unsure{arithmatisation|arithmétisation} de $\pi$
--- et, pour commencer, de retrouver
$\unsure{\mathcal N_a}$ (dans lequel $\Gamma'\cdot\Sigma$ est d'indice
fini) ? \deleted{Est} Il \unsure{suffirait|suffira}, pour pouvoir
répondre par l'affirmative, de savoir
que tt \unsure{isom.} ext. $\pi\longrightarrow\Pi_{g,\nu}$,
qui envoie \inserted{$\Sigma$ sur $\widehat{\mathcal T}_{g,\nu}=\Sigma_{g,\nu}$, et}
\unsure{le} ss-groupe d'indice
fini $\Gamma'\cdot\Sigma$ dans $\mathcal N_a$ dans $\mathcal N_{g,\nu}$,
est compatible avec les \unsure{arithmatisations|arithmétisations}.
Or ceci revient \illegible{1} :

\textcircled{g} (facultatif \unsure{quand même} !)
\deleted{\unsure{tout}} Pour tt sous-groupe
ouvert $\unsure{N'|\mathcal N'}$ de $\mathcal N_{g,\nu}$,
les éléments de $\widehat{\widehat{\mathcal T}}_{g,\nu}$
\inserted{normalisant $\widehat{\mathcal T}_{g,\nu}$ et}
qui transportent
$\unsure{N'|\mathcal N'}$ dans $\mathcal N_{g,\nu}$, \unsure{sont} dans $\mathcal N_{g,\nu}$
(\unsure{a fortiori},
$\mathcal N_{g,\nu}$ serait son propre normalisateur
dans $\operatorname{Norm}_{\widehat{\widehat{\mathcal T}}_{g,\nu}}
(\widehat{\mathcal T}_{g,\nu})$.
\margin{404}

\scanpage{406}
Avec les ``points de $U$'' (définis comme \unsure{classes}
de conjugaison de \unsure{relèvts} (19)
``admissibles'' i.e. satisfaisant (20), on
fait un \unsure{groupoïde}, ayant comme
groupe fondamental extérieur $\pi$, et sur
lequel $\Gamma^{\unsure{\natural|\sharp}}\subset\mathcal N_a$ opère strictement
(\uline{NB} pour se \unsure{reposer|repérer}, on se donne
maintenant $\unsure{\Gamma^{\natural}|\Gamma}$ lui-$\overline{\mathrm m}$, pas \unsure{seulement} son
germe --- \unsure{moralmt|moralement}, cela signifie
qu'on a une courbe définie sur une
sous-ext. finie $K$ de $\overline{\mathbb Q}/\mathbb Q$ \ldots\unsure{\}|)}.
Les pts fixes de $\Gamma$ correspondent
aux pts rat/$K$, les pts fixes sous un
ss-groupe fermé $\Gamma'$ aux pts rationnels
sur la sous-extension $K'$ de $\overline{\mathbb Q}/\mathbb Q$ asso-
ciée à $\unsure{\Gamma'|\Gamma_{\bullet}'}$ \ldots)
\margin{405}

\scanpage{407}
\mspage{206}

Soit maintenant $I'\subset I=I(\pi)$ une
partie de $I$, stable par $\Gamma$.
\figure{[Au-dessus, $I'\subset I=I(\pi)$ est entouré d'un encadrement par le haut et sur les côtés, ouvert en bas.]}
On \unsure{trouve}
par "\unsure{bouchage} des trous en $I'$" un
groupe ext. à lacets $\pi'$, \deleted{\unsure{qui}} et un
hom extérieur
\[
(20)\qquad \pi^{\unsure{\natural|\sharp}}\longrightarrow\pi'
\]
D'ailleurs $\pi'$ hérite d'une arithmétisation \inserted{$\unsure{a'}$}
par $\pi$ -- et \unsure{on a}
\[
(21)\qquad
\left\{
\begin{array}{ccc}
\Gamma_a & \overset{\sim}{\longrightarrow} & \Gamma_{a'} \\
\unsure{P|\mathbb P}_a & \simeq & \unsure{P|\mathbb P}_{a'} \\
\unsure{\wr} & & \\
\unsure{P|\mathbb P} & &
\end{array}
\right.
\quad\unsure{\leadsto}
\]
donc $\pi'$ est défini sur le $\overline{\mathrm m}$ $\unsure{P|\mathbb P}$ que
$\unsure{\Gamma_a|\pi_a}$
(i.e. en faisant des trous dans $U$ pour
\unsure{trouver} $U'$, on n'a pas changé le
corps de base alg. clos $\overline{\mathbb Q}\ldots$). D'ailleurs
on aura un hom canonique
\[
\widehat{\widehat{\mathcal T}}(\pi)\longrightarrow
\widehat{\widehat{\mathcal T}}(\pi')
\]
induisant
\[
(22)\qquad
\]
\diagram{[Rangée supérieure : $\mathcal N_a\longrightarrow\mathcal N_{a'}$.
Rangée inférieure : $\widehat{\mathcal T}_a\longrightarrow\widehat{\mathcal T}_{a'}$.
Entre chaque objet supérieur et l'objet placé au-dessous, un signe vertical
\unsure{d'inclusion en forme de $\cup$, ouvert vers l'objet supérieur|de flèche descendante très courte}.]}

\margin{406}

\scanpage{408}

(en \unsure{donnant|déduisant} \unsure{justement} $\Gamma_a\longrightarrow\Gamma_{a'}$ par passage
aux quotients), et on trouve, en
composant
\[
\Gamma\;\unsure{\subset|\hookrightarrow}\;\mathcal N_a\longrightarrow\mathcal N_{a'}
\]
un hom \unsure{égalmt|également} injectif
\[
(23)\qquad\Gamma\hookrightarrow\mathcal N_{a'}
\]
qui est une \unsure{quasi-section|quasi section} de $\mathcal N_{a'}$ sur $\Gamma_{a'}$,
\unsure{Donc|On} \deleted{par}\inserted{sous} réserve d'admissibilité,
on trouve sur \deleted{\illegible{1}} $\pi$ une structure
de courbe algébrique potentielle, relatif
au $\overline{\mathrm m}$ $\unsure{P|\mathbb P}$.

\deleted{L'admissibilité s'exprime
\unsure{peut-être par} la condition
\[
(24)\qquad {\pi'}^{\Gamma}=(1)
\]
-- et s'il le faut, on mettra ces
conditions (pour \unsure{tte|$\forall$} $I'\subset I$) dans
la définition de l'admissibilité de}

\deleted{\illegible{1}} Je suis vraiment gêné
aux entournures, faute d'avoir une
définition en forme de "admissible" --

\margin{407}

\scanpage{409}
\mspage{207}

je vais y revenir très vite --- mais
pour le moment, j'ai envie de noter que,
\deleted{du fait que $\Gamma$ \unsure{opérait} \unsure{tr.}}
\inserted{si $I'\neq\varnothing$,}
le groupe ext. $\pi'$ peut se décrire par
un \inserted{vrai} groupoïde, ayant $I'$ comme ens.
d'objets, et sur lequel $\Gamma$ opère au
sens strict (ceci est de l'algèbre \unsure{pure},
indépendamment des histoires d'arithmétisation\ldots)
Le fait que $\Gamma$ opère \unsure{trivialt}
sur les $i'\in I'$ implique que pour \unsure{tt|$\forall$}
$i'\in I'$, on a un \deleted{\unsure{relèvement de}}
\unsure{hom.} canonique
\[
(24)\qquad \Gamma\longrightarrow\pi'(i')
\]
qui relève son action extérieure, d'où
une \unsure{uq} classe de \inserted{$\pi'$-}conj. de relèvements
de l'action ext. de $\Gamma$ sur $\pi'$
en une vraie action de $\Gamma$ sur $\pi'$
--- on espère \unsure{qu'elle} \unsure{satisfait} ${\pi'}^{\Gamma}=1$ ---
et on a donc une application can.
\[
(25)\qquad I'\longrightarrow\mathrm{Pts}(\pi',\unsure{\alpha'|a'},\Gamma)
\]

\margin{408}

\scanpage{410}

Je \uline{dis que cette application est injective}

Ceci est "évident" quand on interprète
"action ext. admissible" par "\unsure{réalisable}
par une vraie courbe alg.", et
"relèvements \unsure{admissibles}" par "\unsure{réalisables}
par des vrais pts rat / \unsure{$K^+$|$K$}, corps des
\unsure{fixes} de $\Gamma$". Mais on voudrait
bien sûr des raisons internes :
la donnée des $(N_{g,\nu})$, et des \unsure{)}
propriétés de ces derniers ! \deleted{\unsure{Et}}

\margin{On va revenir au cas \unsure{card} $I'=2$\ldots}

Ce pt étant admis (en mettant
entre parenthèses les deux définitions
\unsure{éventuelles} d'admissibilité, sur lesquelles
on va revenir plus bas) on trouve
un foncteur "bouchage de trous"
\diagram{\begin{tikzcd}[ampersand replacement=\&,column sep=large]
\begin{gathered}
\text{Courbes algébriques}\\
\text{potentielles sur }P,\ \text{\unsure{\deleted{relatives}|relatives}}\\
\text{à un ss-groupe ouvert}\\
\text{\deleted{\unsure{d'un ens}}}\\
\Gamma\text{ de }\unsure{\Gamma_P|\Pi_P},\ \text{et munies d'un}\\
\text{ens }I'\text{ de pts à l'}\infty
\end{gathered}
\arrow[r]
\&
\begin{gathered}
\text{Courbes alg.}\\
\text{pot. sur }P\text{ relatives}\\
\text{à }\Gamma,\ \text{munies}\\
\text{d'une partie}\\
\text{de l'ens des}\\
\text{"pts \unsure{ratls|réels}/}\Gamma\text{"}
\end{gathered}
\end{tikzcd}}

\margin{il faut pour ceci les \unsure{fonctorialités|constructions}
\unsure{anabéliennes} \unsure{habituelles}}

\margin{409}
